\chapter{DNS und DNSSEC}
\label{section-DNSSEC}
%Domain Name System oder kurz DNS ist ein dienst, der es ermöglicht über Datenbanken einen Host Namen in eine IP Adresse aufzulösen und umgekehrt.
%Es ist ein wichtiger Dienst um kommunikation im Internet zu ermöglichen.
%DNS wurde entworfen möglichst leichtgewichtig zu sein, mit kleinen Paketen und wenig overhead (durch UDP), und dadurch auch geringeren Latenzen.
%Das DNS Protokoll läuft standartmäßig auf UDP. TPC-fallback ist möglich wenn die Nachrichten wegen der Größe fehlschlagen sollten, es ist aber auch möglich direkt DNS Pakete über TCP zu senden, mit gesetztem Bit.\cite{ietf-rfc-1034}
%Die UPD Nachrichtengröße ist auf 512 Byte beschränkt, für DNS Nachrichten reicht das eigentlich aus.\cite{ietf-rfc-1034}\cite{ietf-rfc-1035}
%DNS wurde entworfen möglichst leichtgewichtig zu sein, mit kleinen Paketen und wenig overhead (durch UDP), und dadurch auch geringeren Latenzen.
DNS ermöglicht eine effiziente und hierarchisch strukturierte Namensauflösung.
Allerdings bietet DNS standartmäßig kaum Sicherheitsfeatures.
DNS geht immer davon aus, dass alle Antworten vertrauenswürdig sind, dadurch ist es möglich gefälschte DNS-Antworten einzuschleußen und Clients, die diese Adresse abfragen wollen, auf eine beliebige Adresse umzuleiten.
Ein Angriff, mit gefälschten DNS-Antwort die Verkehr umzuleiten, ist DNS-Cache-Poisoning.

\section{DNS-Cache-Poisoning Angriff}
DNS nimmt alle DNS-Antworten entgegen, ohne diese kryptografisch auf Echtheit zu überprüfen.
Dadurch kann ein Angreifer, gefälschte Resource Records an den DNS-Server senden, die anstelle der angegebenen Website, auf eine eine schädliche, vom Angreifer kontrollierten Seite, verweisen.

DNS nimmt diesen Eintrag einfach entgegen und speichert ihn im Cache.
Wenn nun jemand eine Anfrage an den DNS Server nach diesem RR stellt, schickt DNS den manipulierten Eintrag aus dem Cache, unwissentlich, dass dieser gefälscht ist und auf einen schädlichen Server verweist.

Ein DNS DNS-Cache-Poisoning Angriff, läuft folgendermaßen ab:

\begin{description}
    \item[Angreifer schickt einen gefälschten Eintrag] Der Angreifer schickt dem DNS-Server einen vermeintlichen Eintrag, der vorgibt von der zuständigen Zone zu stammen.
    \item[Die Antwort wird angenommen] Da die Herkunft und somit auch die Integrität der Antwort nicht überprüft werden kann, wird sie einfach vom Resolver angenommen.
    \item[RR wird im Cache zwischengespeichert] Damit der Resolver nicht jedes Mal erneut nach einem RR suchen muss, wird dieser zum schnelleren Zugriff, im Cache zwischengespeichert.
    \item[Client fragt nach dem Eintrag] Ein Client fragt nach dem RR. Der Server antwortet mit dem gefälschten RR aus seinem Cache.
    \item[Erfolgreicher Angriff] Das Opfer wurde nun durch den, vom Angreifer gefälschten RR, unwissentlich mit dem Server des Angreifers \cite{asudbringUebersichtUeberDNSSEC}.
\end{description}
Damit wurde ein erfolgreicher Angriff durchgeführt und der Angreifer kann sich somit, einfach als den angegebene Server ausgeben.
Dadurch sind nun weitere Angriffe oder sonstige schädliche Vorhaben möglich.


\section{Die DNSSEC Erweiterungen}
Um die Authentizität und Integrität der DNS Einträge zu gewährleisten, gibt es die DNS Security Extensions (DNSSEC).
DNSSEC signiert Resource Records und gewährt somit Authentizität und Integrität, die von einem Resolver geprüft werden können.
Es muss berücksichtigt werden, dass dadurch aber keine Vertraulichkeit der Daten gewährleistet wird und auch Angriffe wie DNS-Rebinding Angriffe und D(DoS) Angriffe sind weiterhin möglich.


\section{DNSSEC Resource Record Typen}
Um eine PKI für DNS möglich zu machen, müssen zusätzlich neue RR und Schlüssel hinzukommen.
Die wichtigsten sind, DNSKEY, RRSIG sowie DS und NSEC(3).

\subsection{ZSK und KSK}
Für DNSSEC sind zwei neue Schlüssel erforderlich:
\begin{itemize}
    \item \textbf{Zone Signing Key (ZSK)} ein privater Schlüssel um die DNSKEY-Record der Zone zu signieren.
    \item \textbf{Key Signing Key (KSK)} ein privater Schlüssel um den Schlüssel einer Zone (ZSK) zu signieren, dieser hat oftmals eine längere Gültigkeitsdauer, als ein ZSK.
\end{itemize}
Die beiden Schlüssel zum Verifizieren, die öffentlichen Schlüssel, für ZSK und KSK sind im DNSKEY Record enthalten.

\subsection{DNS Public Key (DNSKEY)}
Der öffentliche Schlüssel für die Zone. Mit diesem RR kann der Resolver die digitalen Signaturen der Zone prüfen.
Dieser Eintrag speichert einen Schlüssel und keine Signatur, anderst als RRSIG.
Für untergeordnete Zonen wird dieser Eintrag als child DNSKEY (CDNSKEY) bezeichnet.

\subsection{RRSIG}
Wenn eine Signatur für eine Zone erstellt wird, wird nicht nur ein einzelner RR signiert, sondern alle RR der Zone werden zu einem RRSET zusammengefasst und dieses RRSET wird zusammen mit den RRSIG-Daten signiert.
Die Signatur wird nach RFC-4033, durch die konkatenation der RRSIG Daten und aller RR berechnet:
\begin{align*}
    signature = sign(RRSIG\_RDATA\ |\ RR(1)\ |\ RR(2)\ldots\ )
\end{align*}
Diese Signatur wird dann in einem RRSIG der Zone gespeichert \cite{ietf-rfc-4033,WhatAreRRSETs2026}.

In Abbildung-\ref{fig:rrsig} ist der Output von \textbf{dig example.com RRSIG} zu sehen. Das ist ein RRSIG Eintrag der Domäne example.com.

\begin{figure}[htbp]
    \lstinputlisting{src/rr_rrsig.txt}
    \caption{RRSIG von example.com}
    \label{fig:rrsig}
\end{figure}

\subsection{Delegation Signer (DS)}
Verknüpft eine delegierte Kindzone mit der übergeordneten Zone.
Es ist ein RR der zum Schützen Delegierung verwendet wird.
Es gibt diesen RR auch für untergeordnete Delegierungssignaturgeber, sog. child delegation signer (CDS).
\subsection{Next Secure (NSEC)}
NSEC ist ein Beweis der nicht-existenz einer Zone oder eines Domänennamens.
NSEC3 verhindert durch hashing einen Angriff durch Zonenaufzählungen, der durch NSEC möglich wird
\cite{ietf-rfc-4033,ietf-rfc-4034,ietf-rfc-4035,asudbringUebersichtUeberDNSSEC,HowDoesDNSSEC}.

\section{DNSSEC trust chain}

Wenn der DNS-Namensraum aus \ref{fig:dns-baum} um die DNSSEC Erweiterungen ergänzt wird, dann könnte er aussehen, wie in \ref{fig:dnssec-baum} dargestellt.

\begin{figure}[htbp]
\centering
\begin{forest}
for tree={
  rectangle, draw, rounded corners=2pt, align=center, font=\small,
  minimum height=1.6em, l sep=32pt, s sep=12pt,
  edge={-{Stealth[length=4pt]}}
}
[{. (Root)\\\scriptsize KSK, ZSK}, fill=yellow!20
  [{com\\\scriptsize KSK, ZSK}, edge label={node[midway,left,font=\scriptsize]{DS}}
    [{example.com\\\scriptsize KSK, ZSK}, edge label={node[midway,left,font=\scriptsize]{DS}}
      [{www.example.com\\\scriptsize A, RRSIG}]
      [{mail.example.com\\\scriptsize MX, RRSIG}]
    ]
  ]
  [{de\\\scriptsize KSK, ZSK}, edge label={node[midway,left,font=\scriptsize]{DS}}
    [{example.de\\\scriptsize KSK, ZSK}, edge label={node[midway,left,font=\scriptsize]{DS}}]
  ]
  [{org\\\scriptsize KSK, ZSK}, edge label={node[midway,right,font=\scriptsize]{DS}}
    [{example.org\\\scriptsize KSK, ZSK}, edge label={node[midway,right,font=\scriptsize]{DS}}]
  ]
]
\end{forest}
\caption{DNS-Namensraum aus \ref{fig:dns-baum} mit DNSSEC Erweiterung}
\label{fig:dnssec-baum}
\end{figure}

Jede Zone besitzt einen ZSK und KSK, sie erzeugen ein RRSIG aus ihren RR und RR-Daten.
Der öffentliche Schlüssel, der obersten Zone, ist der \textbf{Vertrauensanker (trust anchor)} des gesamten Namensraumes, von diesem aus, werden alle weiteren Schlüssel verifiziert.
Alle Elternzonen enthalten DS Records, um zu bestätigen, welche Kinderzonen zu ihnen gehören \cite{asudbringUebersichtUeberDNSSEC,ietf-rfc-4033,HowDoesDNSSEC}.

\newpage

\subsection{Signatur Algorithmen für DNSSEC}
Durch diese Erweiterungen von DNSSEC ist eine vergrößerung der UDP Pakete (Extension Mechanisms for DNS (EDNS0)) notwendig \cite{ietf-rfc-6891}.
von DNNSEC werden die Algorithmen aus Abbildung-\ref{fig:DNSSECalgos} derzeit unterstützt \cite{roseSecureDomainName2026}.
\begin{table}[htbp]
\centering
\begin{tabular}{|c|c|>{\centering\arraybackslash}p{3.5cm}|c|}
    \hline
    Signatur Algorithmus & DNSSEC Code & Länge öffentlicher Schlüssel (in Bits) & Signatur Länge (in Bits) \\
    \hline
    RSA mit SHA-256 & 8 & $\ge 2048$ & $\approx2050-4100$ \\
    \hline
    ECDSA P-256 mit SHA-256 & 13 & 256 & 512 \\
    \hline
    ECDSA P-384 mit SHA-384 & 14 & 384 & 768 \\
    \hline
    Ed225516 & 15 & 256 & 512 \\
    \hline
    Ed448 & 16 & 456 & 912 \\
    \hline
\end{tabular}
\caption{von DNSSEC Unterstützte Signatur Algorithmen}
\label{fig:DNSSECalgos}
\end{table}

Wie aus Abbildung-\ref{fig:rrsig} anhand des DNS Code 13 entnehmen kann wird, in diesem Paket, ECDSA P-256 mit SHA-256 als Signatur Algorithmus verwendet.

\subsubsection{Upgrade auf PQC}
Die klassischen Signatur Algorithmen haben bereits große Schlüssel- und Signaturlängen und DNSSEC benötigt eine Erweiterung der UDP Pakete.\cite{ietf-rfc-2671}
Die Post-Quantum Signatur Algorithmen gehen teilweise bereits in den KB Bereich. Zum Vergleich in \ref{fig:NIST-PQC-Standarts}, die NIST-Standarts für Digitale Signaturen und deren Varianten.
\begin{table}[htbp]
    \centering
    \begin{tabular}{|c|c|>{\centering\arraybackslash}p{3.5cm}|c|c|}
        \hline
        PQC- Algorithmus & Typ & Öffentlicher Schlüssel Länge & Signatur Länge\\
        \hline
        CRISTALS-DILITHIUM2 & Gitterbasiert & 1312 Bytes & 2420 Bytes\\
        \hline
        CRISTALS-DILITHIUM3 & Gitterbasiert & 1952 Bytes & 3293 Bytes\\
        \hline
        CRISTALS-DILITHIUM5 & Gitterbasiert & 2592 Bytes & 4595 Bytes\\
        \hline
        FALCON-512 & Gitterbasiert NTRU & 897 Bytes & 666 Bytes\\
        \hline
        FALCON-1024 & Gitterbasiert NTRU & 1793 Bytes & 1280 Bytes\\
        \hline
        SPHINCS$^+$-128s & Hashbasiert & 32 Bytes & 8\,080 Bytes \\
        \hline
        SPHINCS$^+$-128f & Hashbasiert & 32 Bytes & 16\,976 Bytes \\
        \hline
        SPHINCS$^+$-192s & Hashbasiert & 48 Bytes & 17\,064 Bytes \\
        \hline
        SPHINCS$^+$-192f & Hashbasiert & 48 Bytes & 35\,664 Bytes \\
        \hline
        SPHINCS$^+$-256s & Hashbasiert & 64 Bytes & 29\,792 Bytes \\
        \hline
        SPHINCS$^+$-256f & Hashbasiert & 64 Bytes & 49\,216 Bytes \\
        \hline
    \end{tabular}
    \caption{NIST PQC-Signatur Standarts}
    \label{fig:NIST-PQC-Standarts}
\end{table}
%SPHINCS$^+$-SHAKE256-128s/128f/192s/192f/256s/256f-simple/robust & Hashbasiert & 32/48/64 Bytes & 8 080/16 976/17 064/35 664/29 792/49 216 Bytes\\

Wie in Abbildung-\ref{fig:NIST-PQC-Standarts} zu sehen ist, sprengen bereits die kleinsten Signatur- und Schlüssellängen den Rahmen von einfachen UDP Paketen.
Jede der SPHINCS$^+$-Varianten kann mit mit den Hashfunktionen, SHAKE oder SHA betrieben werden \cite{Falcon,PQCRainbow,schwabeDilithium,computersecuritydivision3SubmissionsPostQuantum2017,bernsteinSPHINCSSignatureFramework2019}.